\section{Discussion}
\label{sec:discussion}

P2E treats human feedback as evidence for revising an evaluation procedure.
Its two stages address different bottlenecks: mining determines which judgments
are requested, while ERA determines how those judgments change the program.
A favorable result would need to survive both independent source-level assessment
and cost-matched comparisons. Agreement alone would not establish that the
evaluator's critique helps generation or that its reward remains useful under
policy optimization.

\subsection{Limitations and future work}

First, available feedback may cover only overall preference. Dimensional
supervision and expert diagnostics require additional annotation and construct
validation. Second, disagreement can reflect noise or shared model bias rather
than useful evaluator gaps; random anchors and source-dependence audits must
precede claims of annotation efficiency. Third, program edits interact, so
connected-candidate gains require controlled attribution rather than causal
interpretation from traces alone. Finally, pointwise scoring assumes a useful
within-domain ordering and adaptive development can overfit. Independent
ID/OOD evaluation, explicit tie policies, and generation-level studies are
needed to determine where this formulation is adequate. Empirical improvements,
transfer across evaluation standards, and reward-training benefits remain open
questions.
